
%(BEGIN_QUESTION)
% Copyright 2006, Tony R. Kuphaldt, released under the Creative Commons Attribution License (v 1.0)
% This means you may do almost anything with this work of mine, so long as you give me proper credit

Ortofosforsyre (H$_{3}$PO$_{4}$) dannes ved å kombinere tetrafosfordekoksid (P$_{4}$O$_{10}$) med vann (H$_{2}$O). Skriv en balansert reaksjonsligning som viser alle reaktanter og alle reaksjonsprodukter i riktige proporsjoner.

\vskip 10pt

Bestem også hvor mange {\it mol} tetrafosfordekoksid som må tilsettes tjue mol vann for at alt skal reagere fullstendig, og hvor mange {\it mol} syre som dannes som resultat.

\vskip 20pt \vbox{\hrule \hbox{\strut \vrule{} {\bf Forslag til sokratisk diskusjon} \vrule} \hrule}

\begin{itemize}
\item{} Forklar hvordan du kan kontrollere arbeidet ditt for å sikre at ligningen er korrekt balansert.
\end{itemize}

\underbar{file i00901}
%(END_QUESTION)





%(BEGIN_ANSWER)

\noindent
{\bf Delvis svar:}

\vskip 10pt

3,33 mol $\hbox{P}_4\hbox{O}_{10}$ trengs for å reagere fullstendig med 20 mol vann, fordi forholdet mellom vann og tetrafosfordekoksid i den balanserte ligningen er 6:1.

%(END_ANSWER)





%(BEGIN_NOTES)

Balansering av reaksjonen med samtidige lineære ligninger:

% No blank lines allowed between lines of an \halign structure!
% I use comments (%) instead, so Tex doesn't choke.

$$\vbox{\offinterlineskip
\halign{\strut
\vrule \quad\hfil # \ \hfil & 
\vrule \quad\hfil # \ \hfil & 
\vrule \quad\hfil # \ \hfil & 
\vrule \quad\hfil # \ \hfil \vrule \cr
\noalign{\hrule}
%
% First row
1 & $x$ & = & $y$ \cr
%
\noalign{\hrule}
%
% Another row
P$_{4}$O$_{10}$ & H$_{2}$O & $\to$ & H$_{3}$PO$_{4}$ \cr
%
\noalign{\hrule}
} % End of \halign 
}$$ % End of \vbox

% No blank lines allowed between lines of an \halign structure!
% I use comments (%) instead, so Tex doesn't choke.

$$\vbox{\offinterlineskip
\halign{\strut
\vrule \quad\hfil # \ \hfil & 
\vrule \quad\hfil # \ \hfil \vrule \cr
\noalign{\hrule}
%
% First row
{\bf Element} & {\bf Balanse-ligning} \cr
%
\noalign{\hrule}
%
% Another row
Fosfor & $4 + 0x = 1y$ \cr
%
\noalign{\hrule}
%
% Another row
Hydrogen & $0 + 2x = 3y$ \cr
%
\noalign{\hrule}
%
% Another row
Oksygen & $10 + x = 4y$ \cr
%
\noalign{\hrule}
} % End of \halign 
}$$ % End of \vbox

Fra fosforbalansen ser vi at $y = 4$. Setter vi denne verdien av $y$ inn i hydrogenbalansen ($2x = 12$), får vi $x = 6$. Setter vi denne verdien av $x$ inn i oksygenbalansen ($10 + 6 = 4y$), får vi $y = 4$, som bekrefter første løsning. Derfor:

$$\hbox{P}_4\hbox{O}_{10} + \hbox{6H}_2\hbox{O} \to \hbox{4H}_3\hbox{PO}_4$$

\vskip 10pt

3,33 mol $\hbox{P}_4\hbox{O}_{10}$ trengs for å reagere fullstendig med 20 mol vann, siden den balanserte ligningen gir et molekylforhold vann:dekoksid på 6:1. 20/6 = 3,33.

\vskip 10pt

13,33 mol syre dannes når tjue mol vann reagerer, siden den balanserte ligningen gir forholdet syre:vann = 4:6.  $(20)(4/6) = 13,33$.



%INDEX% Chemistry, stoichiometry: balancing a chemical equation

%(END_NOTES)

